\section{The upper bound by enriched affine comparison}
\label{upper:section}

We continue to assume $f\in C_c^\infty(\mathbb R)$. The upper
bound uses a finite cascade as an additional Gibbs coordinate.
A Gaussian perturbation with vanishing pressure cost forces the joint
overlap identities at deterministic minimizers of a comparison functional.
The field derivatives at those minimizers give precisely the tail
order needed in Proposition~\ref{var:continuity-order}.
The comparison combines Gaussian interpolation in the tradition of
Guerra~\cite{guerra2003} with Mourrat's enriched supersolution
method~\cite[Section~4]{mourrat2021}. In particular, the use of penalized
deterministic contact points to obtain overlap identities follows the
strategy of that work. We give the estimates for the present model,
where Poisson differentiation of the pattern density replaces the
Gaussian interaction derivative.

\begin{proposition}[Upper bound]
\label{upper:bound}
For every nondecreasing $[0,1]$-valued path $q$,
\begin{equation}
\label{upper:conclusion}
 \limsup_{N\to\infty}\mathbb E p_N
 \le \alpha V_f(q)+S_{\mathrm I}(q).
\end{equation}
Consequently the limit superior is bounded above by the
variational infimum in Theorem~\ref{thm:main}.
\end{proposition}

\subsection{The enriched model and its finite-size estimates}

We use the cascade-field enrichment of Mourrat and
Panchenko~\cite{mourratpanchenko2020}: the cascade label is an additional
Gibbs coordinate, and a hierarchical Gaussian magnetic field probes the
spin overlap. The diagonal subtraction below uses the same normalization
principle. The nonlinear patterns act only on the spins.

Fix a step path $q$ on the partition
$0=\zeta_0<\cdots<\zeta_k=1$, with
$q_{k-1}<1$ and $S_{\mathrm I}(q)<\infty$.
Use the depth-$(k-1)$ cascade of Lemma~\ref{arr:cascade}, with
leaf weights $v_\gamma$ and common level $\gamma\wedge\eta$.
When $k=1$ the cascade has its single root leaf of weight one.
The reference probability on pairs $(x,\gamma)$ is
$\nu_N(dx)v_\gamma$. Define
\[
 R(x,x')=\frac{x\cdot x'}N,
 \qquad T(\gamma,\eta)=\frac{\gamma\wedge\eta}{k}.
\]
The array $T$ is positive semidefinite: assign independent
Gaussian increments of variance $1/k$ to each edge of the tree
and sum them along each root-to-leaf path. The covariance of
these sums is $T$. Its diagonal is $(k-1)/k$, including the
value zero for $k=1$.

Let $0\le t\le\alpha$ and let $K$ be Poisson with mean $Nt$.
The $K$ independent Gaussian patterns act only on $x$. For a
field path $0\le h_0\le\cdots\le h_{k-1}$, add the energy
\[
 z(\gamma)\cdot x-\frac N2h_{k-1},
\]
where the $N$ coordinates of $z$ are independent centered
Gaussian fields on the cascade, each with covariance
$h_{\gamma\wedge\eta}$. They are constructed from common root
increments of variance $h_0$ and independent vertex increments
of variances $h_i-h_{i-1}$ thereafter.

Enumerate the monomials
\[
 \kappa_j(R,T)=R^pT^d,
 \qquad p,d\in\{0,1,2,\ldots\},\quad p+d\ge1,
\]
so that every such monomial has a fixed finite index $j$.
The convention is $T^0=1$, also when $T=0$.
There are independent centered Gaussian fields $Y_j(x,\gamma)$
with these covariances. One explicit construction uses a standard
Gaussian tensor contracted with $(x/\sqrt N)^{\otimes p}$
to produce covariance $R^p$ (a scalar Gaussian for $p=0$).
For $d>0$, place independent copies at the vertices and multiply
them by square roots of the increments of $(i/k)^d$; for
$d=0$ use just a common root copy. The resulting covariance is
the indicated product. This also covers identically zero
covariances without exception.

All these fields are independent of the patterns and the field
$z$. Set
\begin{equation}
\label{upper:perturbation}
 e_N=N^{-1/16},\qquad
 \sqrt N e_N\sum_{j=1}^N2^{-j}u_jY_j(x,\gamma),
 \qquad 1\le u_j\le2,
\end{equation}
and add this perturbation to the energy. Denote the random
pressure of the resulting model by $P_N(t,h,u)$ and put
\begin{equation}
\label{upper:F}
 F_N(t,h,u)=-\mathbb E P_N(t,h,u).
\end{equation}
Every parameter is held deterministic in this expectation.
The diagonal and the absolute value of every covariance
$\kappa_j$ are at most one, so the perturbation covariance is
bounded by $CNe_N^2$. Lemma~\ref{arr:calculus} therefore shows
that turning the perturbation on or off changes the expected
pressure by at most $Ce_N^2$, uniformly in these parameters.
At the energy scale, however, the coefficient of each fixed $Y_j$
grows like $\sqrt N e_N=N^{7/16}$; this is the scale against
which its fluctuations will be compared.

At $t=0$ with the perturbation removed, spin integration at a
leaf gives the product $\prod_{a=1}^N\cosh z_a(\gamma)$.
The independent-coordinate form of Lemma~\ref{arr:tilt} and
the diagonal subtraction imply
\begin{equation}
\label{upper:initial-value}
 \mathbb E P_N(0,h,\text{no perturbation})=\ell(h).
\end{equation}
With patterns and perturbations present, one can instead apply
that lemma to the complete spin log integral at each leaf.
Conditional on patterns, count, and common root Gaussian data,
the remaining vertex marks have exactly the conditional
independence required there. Hence the sampled label shape has
the original cascade law. In particular, for every $N,t,h,u$,
\begin{equation}
\label{upper:label-law}
 \mathbb E\langle\mathbf 1_{
        \{\gamma^1\wedge\gamma^2=i\}}\rangle=w_i.
\end{equation}

\begin{lemma}[Uniform concentration on a field cap]
\label{upper:concentration}
For each fixed $H<\infty$ there is a constant $C_H$, also
depending on the fixed partition, $\alpha$, and $f$, such that
\begin{equation}
\label{upper:variance}
 \operatorname{Var}(P_N(t,h,u))\le\frac{C_H}{N}
\end{equation}
whenever $0\le t\le\alpha$, $0\le h_0\le\cdots\le
h_{k-1}\le H$, and $u\in[1,2]^N$.
\end{lemma}

\begin{proof}
Write $L_N=NP_N$. Put the patterns, their count, and all common
root Gaussian variables into the root data. Apply
Lemma~\ref{arr:tilt} to the leaf spin log integral. If $X_0$
is the resulting recursion value at the root, that lemma gives
\[
 \mathbb E[L_N\mid\text{root data}]=X_0,
 \qquad
 \mathbb E[(L_N-X_0)^2\mid\text{root data}]\le C_{\zeta}.
\]
The second assertion can also be seen directly: $L_N-X_0$
is the difference between the logarithms of the transformed
and original unnormalized cascade totals. Each total has the
original conditional law, whose logarithm has finite second
moment depending only on the fixed cascade. The two totals
need not be independent for this bound.

Conditional on $K$, the norm of the coefficient vector of the
common Gaussian variables in the energy at any fixed state is
at most $C_H\sqrt N$. Indeed its squared norm is the sum of
the common variances, bounded by $NH+CNe_N^2$.
The spin log integral has the same Lipschitz constant in these
variables: the difference of two log integrals is bounded by
the supremum of the difference of their energies. Every
subsequent log-exponential transform preserves such a uniform
change bound. Thus $X_0$ has Gaussian variance at most $C_HN$
by Lemma~\ref{arr:variance}. Replacing one of the $K$ patterns
changes $X_0$ by at most $2\|f\|_\infty$, because this bound
holds for every leaf energy and is preserved by each recursion.
The independent-input variance bound therefore adds at most
$C K$. Combining these independent blocks gives
\[
 \operatorname{Var}(X_0\mid K)\le C_HN+C K.
\]
The conditional mean $\mathbb E[X_0\mid K=m]$ changes by at
most $\|f\|_\infty$ between $m$ and $m+1$, by adding one
pattern under a coupling. If $K'$ is an independent copy,
the independent-copy variance formula bounds the variance of
this conditional mean by
$\|f\|_\infty^2\mathbb E(K-K')^2/2
=\|f\|_\infty^2Nt$.
The total variance of $L_N$ is consequently $O_H(N)$.
Division by $N^2$ proves \eqref{upper:variance}.
\end{proof}

\begin{lemma}[Derivatives of the comparison function]
\label{upper:derivatives}
The function $F_N$ is continuous on the parameter region above.
For every nonnegative nondecreasing direction $\lambda$ on
the fixed partition, its right directional derivative in $h$ is
\begin{equation}
\label{upper:field-derivative}
 \partial_{h,\lambda}F_N
 =\frac12\mathbb E\langle
       R_{12}\lambda_{\gamma^1\wedge\gamma^2}\rangle.
\end{equation}
For $t>0$, its time derivative is
\begin{equation}
\label{upper:time-derivative}
 \partial_tF_N
 =-\mathbb E\log\left\langle
    \exp f(\widehat g\cdot x/\sqrt N)\right\rangle,
\end{equation}
where $\widehat g$ is an independent standard Gaussian pattern.
\end{lemma}

\begin{proof}
Increasing $h$ by $a\lambda$ adds an independent Gaussian
energy with covariance
\[
 aN R(x,x')\lambda_{\gamma\wedge\eta}.
\]
It also changes the deterministic subtraction by
$-aN\lambda_{k-1}/2$.
The diagonal covariance term in Lemma~\ref{arr:calculus}
cancels this subtraction, proving
\eqref{upper:field-derivative} after reversing the pressure
sign. Poisson differentiation gives
\[
 \frac{d}{dt}\mathbb E L_N(K)
 =N\mathbb E[L_N(K+1)-L_N(K)],
\]
whose increment is the logarithm in
\eqref{upper:time-derivative}. Its absolute value is bounded
by $\|f\|_\infty$, so differentiating the Poisson series
and passing to its limits are justified. Continuity in $t$
follows also by coupling Poisson counts. Continuity in $h$
follows by coupling the vertex increments and the Gaussian
interpolation estimates, including at zero increments;
continuity in $u$ follows from ordinary Gaussian-coupling
differentiation. The countable label set is handled by the
field-independent finite truncations and moment bounds in
Lemma~\ref{arr:calculus}.
\end{proof}

\subsection{Deterministic minimizers and their overlap identities}

The quadratic penalty and concentration argument at the minimizing
parameters implement the contact-point construction in
Mourrat~\cite[Section~4, in particular equations~(4.19)--(4.28)]{mourrat2021}.
The proof is included because the pattern contribution and its
concentration estimates differ from those of the Gaussian bipartite
Hamiltonian treated there.

Suppose, contrary to the step version of
\eqref{upper:conclusion}, that for some $\eta>0$ and an
infinite sequence of dimensions,
\[
 \mathbb E p_N\ge S_{\mathrm I}(q)+\alpha V_f(q)+\eta.
\]
Couple $\lfloor\alpha N\rfloor$ patterns to a Poisson count
of mean $N\alpha$ by using the same initial rows. The pressure
change in expectation is at most
\[
 \frac{\|f\|_\infty}{N}
 \mathbb E|K-\lfloor\alpha N\rfloor|=O(N^{-1/2}).
\]
Choose $\delta>0$ with $\alpha\delta<\eta/4$. On the
compact parameter region
\[
 0\le t\le\alpha,\quad
 0\le h_0\le\cdots\le h_{k-1}\le H,\quad
 u\in[1,2]^N,
\]
minimize the continuous function
\begin{equation}
\label{upper:objective}
 \begin{split}
 J_N(t,h,u)={}&F_N(t,h,u)
       -\frac12\sum_{i=0}^{k-1}w_iq_ih_i
       +S_{\mathrm I}(q)+t\{V_f(q)+\delta\}\\
       &+\sum_{j=1}^N2^{-j}(u_j-3/2)^2.
 \end{split}
\end{equation}
The fixed field cap $H$ will now be chosen.

\begin{lemma}[The field cap is coercive]
\label{upper:coercive}
For the fixed partition,
\begin{equation}
\label{upper:field-growth}
 \ell(h)\le (\mathbb E|G|)\sqrt{h_0}+C_{\zeta}
                   -\frac12\sum_{i=0}^{k-1}w_ih_i.
\end{equation}
For sufficiently large fixed $H$, every sufficiently large
dimension in the violating sequence has a minimizer of
\eqref{upper:objective} satisfying
$t>0$ and $h_{k-1}<H$.
\end{lemma}

\begin{proof}
Start from $\log\cosh x\le|x|$. For $d>0$,
\[
 \mathcal T_{s,d}(|\cdot|+c)(x)
 \le |x|+c+\frac{ds}{2}+\frac{\log2}{d},
\]
because $e^{d|x+\sqrt sG|}\le
e^{d(x+\sqrt sG)}+e^{-d(x+\sqrt sG)}
\le e^{d|x|}(e^{d\sqrt sG}+e^{-d\sqrt sG})$ after
taking expectations. Propagate this bound through levels
$k-1,\ldots,1$ and then average $|\sqrt{h_0}G|$ at the
root. The variance-dependent constant is
\[
 -\frac{h_{k-1}}2
 +\frac12\sum_{i=1}^{k-1}\zeta_i(h_i-h_{i-1})
 =-\frac12\sum_{i=0}^{k-1}w_ih_i,
\]
which proves \eqref{upper:field-growth}, with
$C_\zeta=\sum_{i=1}^{k-1}(\log2)/\zeta_i$; the sum is
zero if $k=1$.

Adding the patterns changes pressure by at most
$\alpha\|f\|_\infty$ in expectation, and the perturbation
cost is $O(e_N^2)$. Thus
\[
 F_N(t,h,u)\ge-\ell(h)-\alpha\|f\|_\infty-O(e_N^2).
\]
Since $S_{\mathrm I}(q)\ge0$, $|V_f(q)|\le\|f\|_\infty$,
and the quadratic penalty is nonnegative, on $h_{k-1}=H$
this gives the uniform lower bound
\begin{equation}
\label{upper:cap-bound}
 J_N(t,h,u)\ge
 \frac{w_{k-1}(1-q_{k-1})}{2}H
 -(\mathbb E|G|)\sqrt H-C_\zeta
 -2\alpha\|f\|_\infty-O(e_N^2).
\end{equation}
The linear coefficient is strictly positive, so choose $H$
large enough that this lower bound is positive for all large
$N$. At $t=0$, the entropy definition and
\eqref{upper:initial-value} instead give
$J_N(0,h,u)\ge-O(e_N^2)$.
Finally, testing $t=\alpha$, $h=0$, and $u_j=3/2$ gives
$J_N\le-\eta/2$ for all sufficiently large dimensions in
the violating sequence, by the Poisson coupling, the choice
of $\delta$, and the perturbation bound. Hence a minimizing
triple lies on neither of the excluded faces.
\end{proof}

Choose a deterministic minimizing triple $(t_N,h_N,u_N)$
for each such dimension. Comparing it with the same $(t_N,h_N)$
and all $u_j=3/2$ shows that
\begin{equation}
\label{upper:penalty}
 \sum_{j=1}^N2^{-j}(u_{N,j}-3/2)^2=O(e_N^2),
 \qquad u_{N,j}=3/2+O_j(e_N)\quad\text{for fixed }j.
\end{equation}
Indeed any two allowed perturbations have expected pressures
within $O(e_N^2)$ of the unperturbed model. In particular,
each fixed coordinate of $u_N$ is eventually interior.

\begin{lemma}[Joint identities at deterministic minimizers]
\label{upper:gg}
Every subsequential limit of the sampled joint arrays $(R,T)$
at these deterministic minimizing triples satisfies the joint
Ghirlanda--Guerra identities of Definition~\ref{arr:gg}.
\end{lemma}

\begin{proof}
Fix $j$. All derivatives in this proof hold $(t_N,h_N)$ and
the other perturbation coordinates fixed. Abbreviate
$a_{N,j}=\sqrt N e_N2^{-j}$ and
$g(v)=\mathbb E P_N(t_N,h_N,u_1,\ldots,v,\ldots,u_N)$.
Ordinary Gibbs differentiation gives
\begin{align}
\label{upper:u-derivatives}
 \partial_{u_j}P_N&=\frac{e_N2^{-j}}{\sqrt N}
                         \langle Y_j\rangle,\\
 g''(u_{N,j})&=e_N^24^{-j}
   \mathbb E\langle(Y_j-\langle Y_j\rangle)^2\rangle.
\end{align}
The random function $P_N$ is convex in $u_j$, since it is
a log integral of an energy affine in that coordinate.
We first control the Gibbs fluctuation of $Y_j$ around
$\langle Y_j\rangle$ by the curvature at the minimum. We then
control the disorder fluctuation of $\langle Y_j\rangle$ by
finite differences. Both estimates hold at the chosen deterministic
coordinate $u_{N,j}$, without averaging it.

Interior minimality of $J_N$ implies
$g'(u_{N,j})=2\cdot2^{-j}(u_{N,j}-3/2)$ and
$g''(u_{N,j})\le2\cdot2^{-j}$. Consequently
\begin{equation}
\label{upper:thermal}
 \mathbb E\langle|Y_j-\langle Y_j\rangle|\rangle
 \le C_j e_N^{-1}.
\end{equation}
For the disorder fluctuation, use minimality over an interval
around $u_{N,j}$. Cancellation of the linear term gives
\begin{equation}
\label{upper:taylor}
 0\le g(u_{N,j}\pm s)-g(u_{N,j})
           \mp s g'(u_{N,j})\le2^{-j}s^2
\end{equation}
whenever both displaced coordinates lie in $[1,2]$.
The first inequality uses convexity of $g$.

For any differentiable convex random function, its derivative
at a point lies between the backward and forward difference
quotients. Apply this fact to $P_N$ with step $s$, and subtract
$g'(u_{N,j})$. The mean quotient errors are at most $2^{-j}s$
by \eqref{upper:taylor}. Replacing each of the three random
function values by its expectation costs at most
$C_H/(s\sqrt N)$ in $L^1$, by
Lemma~\ref{upper:concentration}. Taking the positive and
negative parts therefore yields
\begin{equation}
\label{upper:derivative-fluctuation}
 \mathbb E|\partial_{u_j}P_N-g'(u_{N,j})|
 \le C_js+\frac{C_H}{s\sqrt N}.
\end{equation}
By \eqref{upper:penalty}, $s=N^{-1/4}$ is eventually allowed.
Then the right side is $O_j(N^{-1/4})=o(e_N^2)$.
Multiplying by $\sqrt N/(e_N2^{-j})$ in
\eqref{upper:u-derivatives}, and adding
\eqref{upper:thermal}, gives
\begin{equation}
\label{upper:Y-concentration}
 \mathbb E\langle|Y_j-\mathbb E\langle Y_j\rangle|\rangle
 =o(\sqrt N e_N).
\end{equation}
Explicitly, the two contributions are
$O_j(N^{1/16})$ and $O_j(N^{5/16})$, whereas
$\sqrt N e_N=N^{7/16}$. All assertions concern one fixed
$j$ at a time; no bound uniform in $j\le N$ is needed.

Let $D$ be a bounded continuous function of the first
$n\times n$ joint overlap block, $n\ge2$. Gaussian
integration by parts in the field $Y_j$ gives
\begin{equation}
\label{upper:ibp}
 \mathbb E\langle Y_j(1)D\rangle
 =a_{N,j}u_{N,j}\,
 \mathbb E\left\langle D\left(
    \sum_{l=1}^n\kappa_j(R_{1l},T_{1l})
       -n\kappa_j(R_{1,n+1},T_{1,n+1})\right)\right\rangle.
\end{equation}
The overlap test has no direct dependence on the Gaussian
coordinates of $Y_j$, so Lemma~\ref{arr:calculus} applies.
By \eqref{upper:Y-concentration}, the left side equals
$\mathbb E\langle Y_j\rangle\mathbb E\langle D\rangle
+o(\sqrt N e_N)$. The single-replica version of
\eqref{upper:ibp} reads
\[
 \mathbb E\langle Y_j\rangle
 =a_{N,j}u_{N,j}
   \{\kappa_j(1,(k-1)/k)
          -\mathbb E\langle\kappa_j(R_{12},T_{12})\rangle\}.
\]
Substitute this expression and cancel the constant diagonal.
After division by $n a_{N,j}u_{N,j}$, one obtains
\begin{align*}
 \mathbb E\langle D\kappa_j(R_{1,n+1},T_{1,n+1})\rangle
 ={}&\frac1n\mathbb E\langle D\rangle
             \mathbb E\langle\kappa_j(R_{12},T_{12})\rangle\\
 &+\frac1n\sum_{l=2}^n
       \mathbb E\langle D\kappa_j(R_{1l},T_{1l})\rangle+o(1).
\end{align*}
Pass to a joint array subsequence, which exists because the
entries range over a compact set. The identity holds for every
fixed monomial and hence, by linearity and uniform polynomial
approximation, for every continuous function of a pair $(R,T)$.
Continuous tests of the finite block determine its measures,
so the same identities hold with every bounded measurable
block test. This is the full conditional-law formulation in
Definition~\ref{arr:gg}.
\end{proof}

\subsection{Ordered supports and the derivative contradiction}

The no-crossing argument below is the monotone-coupling mechanism in
spin-glass synchronization; compare Panchenko~\cite[Section~4, Lemma~2]{panchenko2015multispecies}.
Its input here is the joint identity already proved for the spin and tree
overlaps.

\begin{lemma}[Ordered coupling of the spin and tree overlaps]
\label{upper:synchronization}
In any joint array limit from Lemma~\ref{upper:gg}, the
two-replica law of $(R_{12},T_{12})$ is supported on a set
ordered in both coordinates. If $p$ is the quantile of
$R_{12}$, the conditional law of $R_{12}$ given $T_{12}=i/k$
is the law of $p(U)$ for $U$ uniform on
$(\zeta_i,\zeta_{i+1})$.
\end{lemma}

\begin{proof}
The arrays $R$ and $R+T$ are Gram arrays of constant diagonal,
and their scalar identities follow from the joint identities.
Theorem~\ref{arr:ultrametric} therefore makes both arrays
ultrametric and makes the off-diagonal entries of $R$
nonnegative. They are at most one by its Gram diagonal.
The tree array $T$ is already ultrametric.

Suppose the two-replica support contained two strictly crossed
points. Choose small neighborhoods $A,B$ of these points so
that every $(r,t)\in A$ and $(r',t')\in B$ satisfy
$r<r'$ and $t>t'$. Both neighborhoods have positive
two-replica probability. The joint conditional law with
$n=2$ then gives
\[
 \mathbb P\big((R_{12},T_{12})\in A,
               (R_{13},T_{13})\in B\big)
 \ge\tfrac12\mathbb P((R_{12},T_{12})\in A)
              \mathbb P((R_{12},T_{12})\in B)>0.
\]
On this event, ultrametricity of $R$ forces $R_{23}=R_{12}$:
the inequalities for edges $12$ and $23$ give the two opposite
inequalities once $R_{12}<R_{13}$ is used. Likewise,
ultrametricity of $T$ forces $T_{23}=T_{13}$.
It follows that $R_{23}+T_{23}$ is strictly smaller than both
$R_{12}+T_{12}$ and $R_{13}+T_{13}$, contradicting
ultrametricity of $R+T$. Hence such support points cannot
exist.

By \eqref{upper:label-law}, the successive values $i/k$ of
$T_{12}$ have masses $w_i$. The conditional spin-overlap
supports at these levels are ordered: every value at an
earlier level is at most every value at a later level.
Thus sorting their mixture simply concatenates these
conditional laws in blocks of lengths $w_i$. This is exactly
the claimed quantile description, including any atoms shared
by adjacent conditional supports.
\end{proof}

The ordered conditional laws identify field directional derivatives with
weighted block averages of the spin-overlap quantile. Directions that
select a tail of the partition give the comparisons required by
Proposition~\ref{var:continuity-order}.

\begin{proof}[Proof of Proposition~\ref{upper:bound}]
Continue along the violating sequence and its minimizing
triples. Extract a joint array limit and let $p$ be the
spin-overlap quantile given by
Lemma~\ref{upper:synchronization}. For any nonnegative
nondecreasing step direction $\lambda$ on the fixed
partition, sufficiently small positive displacements of $h_N$
are feasible, since $h_{N,k-1}<H$. Minimality and
\eqref{upper:field-derivative} yield
\[
 \mathbb E\langle R_{12}
                  \lambda_{\gamma^1\wedge\gamma^2}\rangle
 \ge\sum_{i=0}^{k-1}w_iq_i\lambda_i.
\]
Pass to the joint array limit and use its ordered conditional
laws to obtain
\begin{equation}
\label{upper:cone-inequality}
 \int_0^1p(u)\lambda(u)\,du
 \ge\int_0^1q(u)\lambda(u)\,du.
\end{equation}
Choose successively $\lambda_i=\mathbf 1_{\{i\ge j\}}$,
$0\le j\le k-1$. This gives
\[
 \int_{\zeta_j}^1p(u)\,du
 \ge\int_{\zeta_j}^1q(u)\,du
 \qquad(0\le j\le k),
\]
where $j=k$ is equality. The function
$s\mapsto\int_s^1p(u)\,du$ is concave, because its
almost-everywhere derivative is the nonincreasing function
$-p(s)$. On each partition interval the corresponding tail
integral of $q$ is affine. Concavity places the first tail
above the chord joining its endpoint values, which in turn
lies above the second tail. Therefore the tail inequalities
hold for every $s\in[0,1]$. By
Proposition~\ref{var:continuity-order},
\begin{equation}
\label{upper:V-comparison}
 V_f(p)\le V_f(q).
\end{equation}

At every minimizing triple $t_N>0$, a small displacement to
the left in time is feasible, even if $t_N=\alpha$ or
$t_N\downarrow0$. Thus the derivative of $J_N$ with respect
to $t$ at that triple is nonpositive. On the other hand, the
fresh pattern in \eqref{upper:time-derivative} has covariance
$R$ and diagonal one. Lemma~\ref{arr:fresh-pattern}, applied
to the marginal limiting spin-overlap array, gives
\[
 \partial_t J_N(t_N,h_N,u_N)
 \longrightarrow -V_f(p)+V_f(q)+\delta\ge\delta>0.
\]
This contradicts the required sign at the minimizers and
proves \eqref{upper:conclusion} for every step trial with
$q_{k-1}<1$ and finite entropy. Infinite-entropy trials give
no restriction. Taking the infimum over finite step trials
and using Corollary~\ref{var:step-reduction} proves the
claimed bound for the full variational infimum, and hence
for every admissible $q$.
\end{proof}
